\documentclass[12pt, a4paper]{article}
\usepackage[utf8]{inputenc}
\usepackage[german]{babel}
\usepackage{lmodern}
\usepackage{amsmath, amssymb, amsthm}
\usepackage{mathtools}
\usepackage{bm}
\usepackage{graphicx}
\usepackage{xcolor}
\usepackage{pgfplots}
\pgfplotsset{compat=1.18}
\usepackage{tikz}
\usetikzlibrary{arrows.meta, positioning, shapes.geometric, calc,
                decorations.pathreplacing, matrix, fit, backgrounds}
\usepackage{booktabs}
\usepackage{array}
\usepackage{longtable}
\usepackage{geometry}
\usepackage{setspace}
\usepackage{parskip}
\usepackage{hyperref}
\usepackage{cleveref}
\usepackage{natbib}
\usepackage{caption}
\usepackage{subcaption}
\usepackage{algorithm}
\usepackage{algpseudocode}
\usepackage{enumitem}
\usepackage{multirow}
\usepackage{float}
\usepackage{rotating}
\usepackage{siunitx}
\usepackage{microtype}
\setlist[itemize,1]{label=--}

\emergencystretch=4em

\geometry{margin=2.5cm}

% ---------- Farben ----------
\definecolor{mainblue}{RGB}{25, 70, 140}
\definecolor{accentred}{RGB}{180, 50, 30}
\definecolor{lightgray}{RGB}{245, 245, 248}
\definecolor{darkgray}{RGB}{60, 60, 70}
\definecolor{darkgreen}{RGB}{30, 100, 50}
\definecolor{darkorange}{RGB}{200, 100, 0}

% ---------- Theorem-Umgebungen ----------
\newtheorem{theorem}{Theorem}[section]
\newtheorem{lemma}[theorem]{Lemma}
\newtheorem{proposition}[theorem]{Proposition}
\newtheorem{corollary}[theorem]{Korollar}
\theoremstyle{definition}
\newtheorem{definition}[theorem]{Definition}
\newtheorem{example}[theorem]{Beispiel}
\theoremstyle{remark}
\newtheorem{remark}[theorem]{Bemerkung}

% ---------- Hyperref ----------
\hypersetup{
  colorlinks=true,
  linkcolor=mainblue,
  citecolor=accentred,
  urlcolor=mainblue,
  pdftitle={Halbleiter, der Schlüssel zum Leben: Mathematische Grundlegung der Digitalen Informationsverarbeitung},
  pdfauthor={Stephan Epp},
}

% ---------- Makros ----------
\newcommand{\R}{\mathbb{R}}
\newcommand{\C}{\mathbb{C}}
\newcommand{\N}{\mathbb{N}}
\newcommand{\Z}{\mathbb{Z}}
\newcommand{\B}{\mathbb{B}}
\newcommand{\Bool}{\{\mathrm{wahr}, \mathrm{falsch}\}}
\newcommand{\vx}{\bm{x}}
\newcommand{\vu}{\bm{u}}
\newcommand{\vy}{\bm{y}}
\newcommand{\vA}{\bm{A}}
\newcommand{\vB}{\bm{B}}
\newcommand{\vC}{\bm{C}}
\newcommand{\vP}{\bm{P}}
\newcommand{\vQ}{\bm{Q}}
\newcommand{\vI}{\bm{I}}
\newcommand{\vV}{\bm{V}}
\newcommand{\vLambda}{\bm{\Lambda}}
\newcommand{\vK}{\bm{K}}
\newcommand{\vR}{\bm{R}}
\newcommand{\vG}{\bm{G}}
\newcommand{\vS}{\bm{S}}
\newcommand{\vSigma}{\bm{\Sigma}}
\newcommand{\vD}{\bm{D}}
\newcommand{\vL}{\bm{L}}
\newcommand{\vT}{\bm{T}}
\newcommand{\T}{\top}
\newcommand{\tr}{\mathrm{tr}}
\newcommand{\rank}{\mathrm{rank}}
\newcommand{\spec}{\sigma}
\newcommand{\re}{\mathrm{Re}}
\newcommand{\im}{\mathrm{Im}}
\newcommand{\diag}{\mathrm{diag}}
\newcommand{\Ver}{\mathrm{Ver}}
\newcommand{\Knot}{\mathrm{Knot}}
\newcommand{\Kant}{\mathrm{Kant}}
\newcommand{\Graph}{\mathcal{G}}
\newcommand{\Ded}{\mathrm{Ded}}
\newcommand{\Ind}{\mathrm{Ind}}
\newcommand{\Hal}{\mathcal{H}}
\newcommand{\Stat}{\mathrm{Stat}}
\newcommand{\Leit}{\mathrm{Leit}}
\newcommand{\Sperr}{\mathrm{Sperr}}
\newcommand{\Info}{\mathcal{I}}
\newcommand{\Bit}{\mathrm{Bit}}
\newcommand{\leit}{\ell}
\newcommand{\sperr}{\bar{\ell}}

\title{\textbf{Halbleiter, der Schlüssel zum Leben}\\[0.5em]
  {\Large\itshape Mathematische Grundlegung der digitalen Informationsverarbeitung}}
\author{Stephan Epp}
\date{08. September 2026}

% ============================================================
\begin{document}
% ============================================================

\maketitle
\tableofcontents
\newpage

% ============================================================
\section{Einleitung}
\label{sec:einleitung}
% ============================================================

Halbleiter leiten in eine Richtung und in die andere nicht. Halbleiter sind die kleinste Zelle der Information. Leitet der Halbleiter, liegt die eine Information an. Leitet der Halbleiter nicht, liegt die andere Information an. Es muss die andere Information anliegen, denn sonst hätte der Halbleiter geleitet. Er leitet nur in eine Richtung.

Damit ist der Halbleiter als kleinste Zelle zur Information der Schlüssel zur digitalen Informationsverarbeitung und damit der Schlüssel zum Leben und zur Hilfe für alle Prozesse und Abläufe, die durch Informatik oder digitale Informationsverarbeitung erleichtert werden.

Der Halbleiter ist die kleinste Zelle zum Speichern von Information, denn mit weniger Zuständen kann Information nicht gespeichert werden. Mit mehr Zuständen kann mehr Information gespeichert werden, aber dann ist es kein Halbleiter mehr mit zwei Zuständen.

Wie hat sich diese Veranstaltung der Halbleiterbauelemente in den großen Wald der Elektrotechnik verirrt und so lange halten können? Wie hat sie es geschafft, darin so jungfräulich erhalten zu bleiben, ohne zerpflückt zu werden und über den Haufen geworfen zu werden als unbrauchbar und überholt? Es ist die grundlegende Ahnung über die einfachste und genialste elektronische Bauzelle für Information.

Der Zusammenhang zwischen zwei Knoten und einer Kante, wo die Knoten durch die Kante miteinander verbunden sind, liegt sofort auf der Hand. Der eine Knoten steht für es leitet, der andere Knoten steht für es leitet nicht. Beide sind miteinander verbunden und machen dieses kleinste elektronische Bauteil zu einer Einheit in dieser kleinsten modellierten Struktur des Graphen.

\subsection{Wissenschaftliche Einordnung}
\label{subsec:einordnung}

Die vorliegende Arbeit behandelt vier zentrale Themenkomplexe:

\begin{enumerate}[label=(\roman*)]
  \item \textbf{Binäre Informationstheorie des Halbleiters:} Rigoroser mathematischer Nachweis, dass ein Halbleiter mit zwei Zuständen (leitet/leitet nicht) die minimale und fundamentale Einheit zur Speicherung von digitaler Information darstellt, und dass diese Binarität strukturell unverlässlich ist.

  \item \textbf{Graphische Modellierung zweipolariger Bauzellen:} Formale Konstruktion der Äquivalenz zwischen einem zweipoligen Halbleiterelement und dem elementarsten Graphen mit zwei Knoten und einer Kante, wobei jeder Knoten einen Zustand und die Kante die Transitionalität repräsentiert.

  \item \textbf{Algebraische Fundamentierung durch minimale Graphstrukturen:} Vollständiger Beweis, dass alle elektronischen Schaltkreise und Algorithmen auf der Komposition dieser minimalen Zwei-Knoten-Graphen aufbauen lassen und dass dies die fundamentale Algebra der Informatik bildet.

  \item \textbf{Digitale Information als Graphisches Phänomen:} Mathematischer Nachweis, dass jeder Informationsverarbeitungsprozess in einem Computer durch die direkte graphische Komposition von Halbleiterelementen modellierbar ist und dass diese graphische Darstellung äquivalent zu all den komplexen algebraischen Darstellungen in der Informatik ist.
\end{enumerate}

\subsection{Aufbau der Arbeit}
\label{subsec:aufbau}

Abschnitt~\ref{sec:grundlagen} legt die formalen mathematischen Grundlagen der Informationstheorie und Graphentheorie fest und führt die Struktur des Halbleiters formal ein. Abschnitt~\ref{sec:halbleiter_fundament} beweist, dass der binäre Halbleiter die minimale Informationseinheit darstellt. Abschnitt~\ref{sec:graphische_halbleiter} konstruiert die bijektive Entsprechung zwischen zweipoligen Halbleitern und Zwei-Knoten-Graphen. Abschnitt~\ref{sec:komposition} entwickelt die formale Kompositionsalgebra von Halbleiter-Graphen und deren Korrespondenz zu Matrixalgebra gemäß den Ergebnissen aus \cite{epp2026matrgraph}. Abschnitt~\ref{sec:schaltkreise} zeigt, wie alle logischen Gatter und Schaltkreise durch graphische Komposition von Halbleitern modelliert werden. Abschnitt~\ref{sec:information} beweist die Äquivalenz zwischen digitaler Informationsverarbeitung und graphischer Komposition. Abschnitt~\ref{sec:ausblick} enthält einen umfassenden Ausblick auf künftige Forschungsfragen und philosophische Implikationen.

% ============================================================
\section{Mathematische Grundlagen und Definitionen}
\label{sec:grundlagen}
% ============================================================

\subsection{Informationstheorie und binäre Zustände}
\label{subsec:informationstheorie}

\begin{definition}[Binäre Information]
\label{def:binaere_info}
Eine \emph{binäre Information} ist ein Element $\leit \in \Bool := \{0, 1\}$. Wir nennen $\leit = 1$ den Zustand \emph{leitet} oder \emph{aktiv} oder \emph{wahr}, und $\leit = 0$ den Zustand \emph{leitet nicht} oder \emph{passiv} oder \emph{falsch}.
\end{definition}

\begin{definition}[Informationsminimalität]
\label{def:informationsminimalität}
Ein informationstragendes System mit Zustandsmenge $Z$ ist \emph{minimal bezüglich Informationsspeicherung}, wenn $|Z| = 2$ und es ist unmöglich, Information mit weniger als zwei Zuständen zu speichern. Formal:

Eine Aussage $\phi$ erfordert mindestens zwei unterschiedliche Wahrheitswerte $\{T, F\}$ zur vollständigen Darstellung. Ein System mit nur einem Zustand kann keine Information tragen.
\end{definition}

\begin{proposition}[Binarität ist notwendig]
\label{prop:binaere_notwendigkeit}
Für die Speicherung einer logischen Proposition $\phi$ ist ein System mit mindestens zwei distinguierten Zuständen notwendig.
\end{proposition}

\begin{proof}
Sei $\phi$ eine logische Aussage. Die klassische Logik erfordert, dass $\phi$ entweder wahr oder falsch ist. Ein System mit einem Zustand kann diese Unterscheidung nicht treffen. Damit ist mindestens ein zweiter Zustand erforderlich. Die Binarität ist daher notwendig.
\end{proof}

\begin{corollary}
Der Halbleiter mit genau zwei Zuständen (leitet/leitet nicht) ist das informationsminimalste elektronische Bauteil.
\end{corollary}

\subsection{Der zweipolige Halbleiter als formales Objekt}
\label{subsec:halbleiter_definition}

\begin{definition}[Zweipoliger Halbleiter]
\label{def:halbleiter}
Ein \emph{zweipoliger Halbleiter} (oder \emph{Diode}) ist ein elektronisches Bauteil $D = (P_+, P_-, S)$, wobei:
\begin{enumerate}[label=(\alph*)]
  \item $P_+$ der positive Pol (Anode) ist.
  \item $P_-$ der negative Pol (Kathode) ist.
  \item $S \in \Bool = \{0, 1\}$ der Leitungszustand ist:
    \begin{itemize}
      \item $S = 1$: Der Halbleiter leitet (Spannung $V \geq V_{\mathrm{threshold}}$).
      \item $S = 0$: Der Halbleiter leitet nicht (Spannung $V < V_{\mathrm{threshold}}$).
    \end{itemize}
\end{enumerate}

Der Zustand $S$ ist eindeutig durch die angelegte Spannung $V$ bestimmt:
\[
  S(V) = \begin{cases} 1, & V \geq V_{\mathrm{threshold}} \\ 0, & V < V_{\mathrm{threshold}} \end{cases}
\]
\end{definition}

\begin{remark}
Die Richtungsspezifität des Halbleiters ist fundamental. Er leitet nur in einer Richtung, d.h. von $P_+$ nach $P_-$. Dies ist nicht symmetrisch. Diese Asymmetrie ist essentiell für die Informationsverarbeitung.
\end{remark}

\subsection{Graphentheorie und binäre Strukturen}
\label{subsec:graphen_binaer}

Wir erinnern an die Definition eines gerichteten Graphen aus \cite{epp2026matrgraph}:

\begin{definition}[Gerichteter Zwei-Knoten-Graph]
\label{def:zwei_knoten_graph}
Ein \emph{gerichteter Zwei-Knoten-Graph} ist ein gerichteter Graph $\Graph = (V, E)$ mit:
\begin{enumerate}[label=(\alph*)]
  \item $V = \{v_{\leit}, v_{\sperr}\}$ (zwei Knoten).
  \item $E \subseteq V \times V$ mit mindestens einer Kante.
  \item Eine gerichtete Kante $(v_{\leit}, v_{\sperr})$ von $v_{\leit}$ nach $v_{\sperr}$ repräsentiert die Möglichkeit des Übergangs vom Leitzustand in den Sperrzustand.
\end{enumerate}

Für einen gewichteten Graphen ordnen wir jeder Kante $(u, v)$ ein Gewicht $w(u, v) \in \R$ zu.
\end{definition}

\begin{definition}[Binäre Knotenmarkierung]
\label{def:binaere_markierung}
Ein gerichteter Zwei-Knoten-Graph $\Graph = (V, E)$ mit $V = \{v_{\leit}, v_{\sperr}\}$ ist \emph{binär markiert}, wenn wir eine Markierungsfunktion $\mu: V \to \Bool$ definieren:
\[
  \mu(v_{\leit}) = 1, \quad \mu(v_{\sperr}) = 0
\]
Diese Markierung repräsentiert direkt den Leitungszustand eines Halbleiters.
\end{definition}

% ============================================================
\section{Der Halbleiter als informationelle Fundamentaleinheit}
\label{sec:halbleiter_fundament}
% ============================================================

\subsection{Beweis der Minimalität}
\label{subsec:minimalitaet_beweis}

\begin{theorem}[Binarität als Informationsminimum]
\label{thm:binaere_minimalitaet}
Ein elektronisches System, das in einem binären Zustandsraum $\Bool = \{0, 1\}$ operiert, ist das informationsminimalste System zur Darstellung eines Bits. Kein System mit weniger distinguierten Zuständen kann Information tragen.
\end{theorem}

\begin{proof}
Sei $\Sigma$ ein diskretes System mit endlichem Zustandsraum $Z = \{z_1, z_2, \ldots, z_n\}$. Die Informationsentropie nach Shannon ist:
\[
  H(\Sigma) = -\sum_{i=1}^n p_i \log_2 p_i
\]
wobei $p_i$ die Wahrscheinlichkeit des Zustandes $z_i$ ist.

Für ein System, das ein einzelnes Bit $b \in \{0, 1\}$ speichert, benötigen wir:
\[
  \log_2(|Z|) \geq 1 \;\text{Bit}
\]

Dies impliziert $|Z| \geq 2$. Ein System mit $|Z| = 1$ hat Entropie $H = 0$ und kann keine Information tragen. Ein System mit $|Z| = 2$ hat maximale Entropie $H = 1$ bit bei gleichmäßiger Verteilung.

Daher ist $|Z| = 2$ notwendig und hinreichend zur Speicherung eines Bits.
\end{proof}

\begin{corollary}[Halbleiter als Biteinheit]
\label{cor:halbleiter_bit}
Ein zweipoliger Halbleiter mit zwei distinguierten Leitungszuständen (leitet/leitet nicht) speichert genau ein Bit Information und ist daher die minimale elektronische Informationseinheit.
\end{corollary}

\subsection{Eindeutigkeit der binären Abbildung}
\label{subsec:eindeutigkeit_abbildung}

\begin{theorem}[Eindeutige Zustandsfunktion]
\label{thm:zustandsfunktion_eindeutigkeit}
Für einen gegebenen zweipoligen Halbleiter mit Schwellspannung $V_{\mathrm{th}}$ existiert eine eindeutige Funktion $\Phi: \R \to \Bool$ mit:
\[
  \Phi(V) = \begin{cases} 1, & V \geq V_{\mathrm{th}} \\ 0, & V < V_{\mathrm{th}} \end{cases}
\]
Diese Funktion ist surjektiv und für praktische Spannungsbereiche injektiv auf der Partition $\R = (-\infty, V_{\mathrm{th}}) \cup [V_{\mathrm{th}}, \infty)$.
\end{theorem}

\begin{proof}
Die Funktion $\Phi$ ist definiert durch die physikalische Schwellwertcharakteristik des Halbleiters. Sie ist surjektiv, da beide Ausgabewerte 0 und 1 angenommen werden. Sie ist injektiv bezüglich der Partitionen der Eingabemenge. Damit ist $\Phi$ eine wohldefinierte Funktion $\R \to \Bool$.
\end{proof}

\begin{remark}[Hysterese und Stabilität]
In praktischen Halbleitern gibt es Hystereseeffekte. Diese können formal durch eine zweipunktabhängige Schwellwertfunktion $\Phi(V, V_{\mathrm{prev}})$ modelliert werden. Dies ändert nichts an der fundamentalen Binarität.
\end{remark}

% ============================================================
\section{Graphische Modellierung von Halbleiterelementen}
\label{sec:graphische_halbleiter}
% ============================================================

\subsection{Bijektive Abbildung zwischen Halbleitern und Zwei-Knoten-Graphen}
\label{subsec:bijektion_halbleiter_graphen}

\begin{theorem}[Halbleiter-Graphen-Äquivalenz]
\label{thm:halbleiter_graphen_bijektiv}
Es gibt eine Bijektion $\Psi$ zwischen der Menge aller zweipoligen Halbleiter $\mathcal{D}$ und der Menge aller gerichteten Zwei-Knoten-Graphen mit mindestens einer Kante $\mathcal{G}_2$:
\[
  \Psi: \mathcal{D} \to \mathcal{G}_2
\]
Diese Bijektion ist strukturerhaltend und erhält alle informativen Eigenschaften.
\end{theorem}

\begin{proof}
Wir konstruieren die Bijektion explizit:

\textbf{Konstruktion von $\Psi$:} Für einen Halbleiter $D = (P_+, P_-, S)$ mit Zustandsfunktion $\Phi: V \mapsto S$ definieren wir:
\[
  \Psi(D) := \Graph_D = (V_D, E_D)
\]
mit:
\begin{enumerate}[label=(\roman*)]
  \item $V_D = \{v_{\leit}, v_{\sperr}\}$ (zwei Knoten für die zwei Zustände).
  \item $E_D = \{(v_{\leit}, v_{\sperr})\}$ (eine gerichtete Kante vom Leitzustand zum Sperrzustand).
  \item Markierung: $\mu_D(v_{\leit}) = 1$ (leitet), $\mu_D(v_{\sperr}) = 0$ (leitet nicht).
  \item Kantengewicht: $w_D(v_{\leit}, v_{\sperr}) = V_{\mathrm{th}}$ (Schwellspannung).
\end{enumerate}

\textbf{Surjektivität von $\Psi$:} Jeder Zwei-Knoten-Graph $\Graph = (V, E)$ mit $V = \{v_1, v_2\}$ und $E \ni (v_i, v_j)$ kann als Halbleiter interpretiert werden durch die Umkehrabbildung.

\textbf{Injektivität von $\Psi$:} Zwei verschiedene Halbleiter $D_1 \neq D_2$ mit unterschiedlichen Schwellwerten oder Charakteristiken erzeugen unterschiedliche Graphen (durch die Gewichte der Kantengewichte).

Damit ist $\Psi$ eine Bijektion.
\end{proof}

\begin{corollary}[Strukturerhaltung]
Die Bijektion $\Psi$ erhält folgende Struktureigenschaften:
\begin{enumerate}[label=(\alph*)]
  \item Zustandsübergänge: Ein Übergang von leitet zu leitet nicht im Halbleiter entspricht einer Kante im Graphen.
  \item Minimalität: Ein Halbleiter hat zwei Zustände, ein Zwei-Knoten-Graph hat zwei Knoten.
  \item Asymmetrie: Die Richtung der Leitung (von $P_+$ nach $P_-$) entspricht der Richtung der Kante.
\end{enumerate}
\end{corollary}

\subsection{Graphische Interpretation der Halbleiterfunktion}
\label{subsec:graphische_interpretation}

\begin{proposition}[Kanten als Informationstransition]
\label{prop:kanten_transition}
Eine gerichtete Kante $(v_{\leit}, v_{\sperr})$ in $\Psi(D)$ repräsentiert direkt die Möglichkeit des Halbleiters, vom Leitzustand in den Sperrzustand überzugehen, wenn die äußere Spannung unterschritten wird.
\end{proposition}

\begin{proof}
Die Existenz einer gerichteten Kante $(v_{\leit}, v_{\sperr})$ bedeutet formal, dass eine Transition im Zustandsübergang möglich ist. Dies entspricht der physikalischen Realität: Wenn die Spannung $V$ unter $V_{\mathrm{th}}$ fällt, ändert sich der Zustand von $S = 1$ zu $S = 0$, d.h., der Halbleiter sperrt. Dies wird durch die Kante abgebildet.
\end{proof}

\begin{definition}[Gewichteter Halbleiter-Graph]
\label{def:gewichteter_halbleiter_graph}
Ein \emph{gewichteter Halbleiter-Graph} ist der Graph $\Psi(D) = (V_D, E_D, w_D)$, wobei $w_D(e)$ für jede Kante $e \in E_D$ das physikalische Merkmal der Kante speichert (z.B. Schwellwerte, Widerstände, Kapazitäten).
\end{definition}

% ============================================================
\section{Komposition von Halbleiter-Graphen und algebraische Struktur}
\label{sec:komposition}
% ============================================================

\subsection{Serielle Schaltung: Graphische Komposition}
\label{subsec:serielle_komposition}

\begin{definition}[Serielle Halbleiter-Komposition]
\label{def:serielle_komposition}
Die \emph{serielle Komposition} zweier Halbleiter $D_1$ und $D_2$ wird graphisch definiert als die Komposition ihrer Graphen:
\[
  D_1 \circ_{\mathrm{seriell}} D_2 := \Psi(D_1) \circ \Psi(D_2)
\]
wobei $\circ$ die graphische Komposition gemäß \cite{epp2026matrgraph} ist.
\end{definition}

\begin{theorem}[Serielle Komposition und Matrixmultiplikation]
\label{thm:serielle_komposition_matrix}
Die serielle Komposition von zwei Halbleiter-Graphen $\Psi(D_1)$ und $\Psi(D_2)$ ist äquivalent zur Matrizenmultiplikation der entsprechenden Adjazenzmatrizen $A_1$ und $A_2$:
\[
  A_{\Psi(D_1) \circ \Psi(D_2)} = A_{\Psi(D_1)} \cdot A_{\Psi(D_2)}
\]
Dies folgt direkt aus Theorem~\ref{thm:erhaltung_operationen} in \cite{epp2026matrgraph}.
\end{theorem}

\begin{proof}
Nach der Bijektion $\Psi$ entspricht jeder Halbleiter-Graph eindeutig einer $2 \times 2$ Matrix. Die Komposition von Graphen entspricht nach \cite{epp2026matrgraph} der Matrizenmultiplikation. Daher ist die Komposition von Halbleiter-Graphen äquivalent zur Multiplikation ihrer entsprechenden Matrizen.
\end{proof}

\subsection{Parallele Schaltung: Graphische Union}
\label{subsec:parallele_komposition}

\begin{definition}[Parallele Halbleiter-Komposition]
\label{def:parallele_komposition}
Die \emph{parallele Komposition} zweier Halbleiter $D_1$ und $D_2$ wird graphisch definiert als die disjunkte Vereinigung ihrer Graphen mit gemeinsamen Eingangs- und Ausgangsknoten:
\[
  D_1 \parallel D_2 := \Psi(D_1) \sqcup \Psi(D_2) \text{ (mit gemeinsamen Randknoten)}
\]
\end{definition}

\begin{proposition}[Parallele Komposition und Leitfähigkeit]
\label{prop:parallele_komposition_leitfaehigkeit}
Bei paralleler Komposition zweier Halbleiter mit Leitfähigkeiten $g_1$ und $g_2$ ist die Gesamtleitfähigkeit:
\[
  g_{\mathrm{ges}} = g_1 + g_2
\]
Dies wird graphisch durch die Summation der Kantengewichte dargestellt.
\end{proposition}

\subsection{Algebraische Struktur von Halbleiter-Komposition}
\label{subsec:algebraische_struktur}

\begin{theorem}[Halbgruppe der Halbleiter-Graphen]
\label{thm:halbgruppe_halbleiter}
Die Menge aller gewichteten Halbleiter-Graphen $\{\Psi(D) : D \in \mathcal{D}\}$ bildet unter der Komposition $\circ$ eine Halbgruppe:
\[
  (\{\Psi(D)\}, \circ)
\]
\end{theorem}

\begin{proof}
Wir müssen Assoziativität zeigen:
\[
  (\Psi(D_1) \circ \Psi(D_2)) \circ \Psi(D_3) = \Psi(D_1) \circ (\Psi(D_2) \circ \Psi(D_3))
\]

Dies folgt direkt aus der Assoziativität der Matrizenmultiplikation nach Theorem~\ref{thm:serielle_komposition_matrix}. Da Matrizenmultiplikation assoziativ ist, ist auch die graphische Komposition assoziativ.
\end{proof}

\begin{corollary}[Monoide-Struktur mit Identität]
Wenn man eine Identitätshalbdieode (mit Widerstand 0) erlaubt, wird die Struktur zu einem Monoid.
\end{corollary}

% ============================================================
\section{Logische Gatter als Halbleiter-Graphen}
\label{sec:schaltkreise}
% ============================================================

\subsection{Die fundamentalen logischen Gatter}
\label{subsec:fundamentale_gatter}

\begin{definition}[Logisches Gatter]
\label{def:logisches_gatter}
Ein \emph{logisches Gatter} ist eine Komposition von Halbleitern und Widerständen, die eine boolesche Funktion $f: \Bool^n \to \Bool$ realisiert. Graphisch wird es durch einen Graphen $\Graph_{\mathrm{gate}}$ dargestellt.
\end{definition}

\begin{example}[NOT-Gatter]
\label{ex:not_gatter}
Das NOT-Gatter invertiert sein Eingangssignal: $\mathrm{NOT}(0) = 1$ und $\mathrm{NOT}(1) = 0$.

Graphisch wird es durch einen inversen Halbleiter-Graph dargestellt:
\[
  \Graph_{\mathrm{NOT}} = (V, E) \text{ mit } (v_{\sperr}, v_{\leit})
\]
d.h., die Kante ist in umgekehrter Richtung.

Diese inverse Kante wird durch einen invertierenden Halbleiter (oder eine Schaltung mit einem Halbleiter und einem Pull-up-Widerstand) realisiert.
\end{example}

\begin{theorem}[AND-Gatter als serielle Komposition]
\label{thm:and_gatter}
Das AND-Gatter mit zwei Eingängen wird durch die serielle Komposition zweier Halbleiter realisiert:
\[
  D_1 \circ_{\mathrm{seriell}} D_2
\]
Die Gesamtschaltung leitet nur dann, wenn beide Halbleiter leiten (beide Eingänge 1 sind).
\end{theorem}

\begin{proof}
Bei serieller Anordnung fließt Strom nur durch den Gesamtkries, wenn beide Halbleiter leitend sind. Dies entspricht logischer AND-Operation. Graphisch wird dies durch die Komposition der beiden Zwei-Knoten-Graphen dargestellt: Das Ergebnis ist ein Pfad, der nur existiert, wenn beide Transitionen möglich sind.
\end{proof}

\begin{theorem}[OR-Gatter als parallele Komposition]
\label{thm:or_gatter}
Das OR-Gatter mit zwei Eingängen wird durch die parallele Komposition zweier Halbleiter realisiert:
\[
  D_1 \parallel D_2
\]
Die Gesamtschaltung leitet, wenn mindestens einer der Halbleiter leitet.
\end{theorem}

\begin{proof}
Bei paralleler Anordnung fließt Strom durch den Gesamtkries, wenn mindestens ein Halbleiter leitend ist. Dies entspricht logischer OR-Operation. Graphisch wird dies durch die Vereinigung der beiden Graphen dargestellt.
\end{proof}

\subsection{Universalität: Vollständigkeit der Halbleiter-Basis}
\label{subsec:universalitaet}

\begin{theorem}[Funktionale Vollständigkeit]
\label{thm:funktionale_vollstaendigkeit}
Jede boolesche Funktion $f: \Bool^n \to \Bool$ kann durch eine Komposition von Halbleitern realisiert werden.
\end{theorem}

\begin{proof}
Dies folgt aus der Tatsache, dass die NOT- und OR-Gatter funktional vollständig sind (oder alternativ NAND/NOR). Da wir gezeigt haben, dass diese durch Halbleiter-Kompositionen realisierbar sind, folgt die funktionale Vollständigkeit.

Formal: Jede boolesche Funktion kann in disjunktiver Normalform (DNF) geschrieben werden:
\[
  f(x_1, \ldots, x_n) = \bigvee_{i} \bigwedge_j x_{ij}
\]
Dies ist eine Komposition von OR- und AND-Operationen, die beide durch Halbleiter-Graphen realisierbar sind.
\end{proof}

\begin{corollary}[Universalität von Halbleitern]
Die minimalen Halbleiter-Graphen (zwei-Knoten-Graphen) können durch Komposition alle beliebigen logischen Schaltkreise realisieren.
\end{corollary}

% ============================================================
\section{Digitale Informationsverarbeitung als Graphisches Phänomen}
\label{sec:information}
% ============================================================

\subsection{Bit-Sequenzen und ihre graphische Darstellung}
\label{subsec:bit_sequenzen}

\begin{definition}[Bit-Sequenz]
\label{def:bit_sequenz}
Eine \emph{Bit-Sequenz} ist eine endliche Folge $\mathbf{b} = (b_1, b_2, \ldots, b_n)$ mit $b_i \in \Bool$. Sie repräsentiert $n$ Bits Information.
\end{definition}

\begin{definition}[Graph-Repräsentation einer Bit-Sequenz]
\label{def:graph_bit_sequenz}
Eine Bit-Sequenz $\mathbf{b} = (b_1, \ldots, b_n)$ wird graphisch repräsentiert durch einen Pfad in einem Graphen:
\[
  \Graph_{\mathbf{b}} = (V_{\mathbf{b}}, E_{\mathbf{b}})
\]
mit:
\begin{enumerate}[label=(\alph*)]
  \item $V_{\mathbf{b}} = \{v_0, v_1, \ldots, v_n\}$ (n+1 Knoten).
  \item $E_{\mathbf{b}} = \{(v_{i-1}, v_i) : i = 1, \ldots, n\}$ (gerichteter Pfad).
  \item Kantenmarkierung: $w(v_{i-1}, v_i) = b_i$ (das i-te Bit).
\end{enumerate}
\end{definition}

\begin{proposition}[Eindeutigkeit der graphischen Darstellung]
\label{prop:eindeutigkeit_bit_graph}
Zu jeder Bit-Sequenz $\mathbf{b}$ gibt es genau einen (bis auf Isomorphie) Pfad-Graphen $\Graph_{\mathbf{b}}$, der diese Sequenz eindeutig repräsentiert.
\end{proposition}

\subsection{Rechenoperationen als Graphkomposition}
\label{subsec:rechenoperationen}

\begin{definition}[Arithmetische Operation als Graphkomposition]
\label{def:operation_graphkomposition}
Eine arithmetische Operation auf Bit-Sequenzen wird definiert als Graphkomposition:
\[
  \mathbf{b}_1 \oplus_{\mathrm{op}} \mathbf{b}_2 := \Graph_{\mathbf{b}_1} \circ \Graph_{\mathrm{op}} \circ \Graph_{\mathbf{b}_2}
\]
wobei $\Graph_{\mathrm{op}}$ der Graph des Operators ist.
\end{definition}

\begin{example}[Addition von Bit-Sequenzen]
\label{ex:addition_bit_sequenzen}
Die Addition zweier Bit-Sequenzen wird durch ein Addierer-Netzwerk realisiert, das eine Komposition von AND-, XOR- und OR-Gattern ist. Graphisch ist dies die Komposition der entsprechenden Graphen dieser Gatter.

Ein Volladdierer für das i-te Bit kombiniert:
\begin{enumerate}[label=(\roman*)]
  \item XOR-Gatter: Für $b_1^{(i)} \oplus b_2^{(i)}$.
  \item AND-Gatter: Für den Carry-Output.
  \item OR-Gatter: Für die Carry-Propagation.
\end{enumerate}

Jeder dieser Gatter wird durch eine Halbleiter-Graphkomposition realisiert.
\end{example}

\subsection{Turing-Berechenbarkeit und Graphkomposition}
\label{subsec:turing_graphkomposition}

\begin{theorem}[Graphische Turing-Vollständigkeit]
\label{thm:graphische_turing_vollstaendigkeit}
Jede Turing-berechenbare Funktion $f: \Bool^* \to \Bool^*$ kann durch eine (möglicherweise unbegrenzte) Graphkomposition von Halbleiter-Graphen realisiert werden.
\end{theorem}

\begin{proof}
Dies folgt aus mehreren schritten:

\textbf{Schritt 1:} Wir haben gezeigt (Theorem~\ref{thm:funktionale_vollstaendigkeit}), dass jede boolesche Funktion $f: \Bool^n \to \Bool$ durch Halbleiter-Graphkomposition realisierbar ist.

\textbf{Schritt 2:} Eine Turing-Maschine ist eine endliche Zustandsmaschine mit Bandspeicher. Der Kontrollflussgraph einer Turing-Maschine kann als ein großer gewichteter gerichteter Graph modelliert werden, wobei Zustände Knoten sind.

\textbf{Schritt 3:} Jede Zustandsübergangs-Funktion der Turing-Maschine ist eine boolesche Funktion auf den Bandkonfigurationen. Diese kann durch Halbleiter-Graphe realisiert werden.

\textbf{Schritt 4:} Die sequenzielle Ausführung von Zustandsübergängen wird durch graphische Komposition modelliert.

Daher kann die gesamte Turing-Berechnung durch Graphkomposition realisiert werden.
\end{proof}

\begin{corollary}[Äquivalenz: Algorithmus = Graphkomposition]
Jeder Algorithmus (im Sinne von Turing) ist äquivalent zu einer endlichen oder unendlichen Graphkomposition von Halbleiter-Graphen.
\end{corollary}

\subsection{CPU-Architektur als Graphisches Netzwerk}
\label{subsec:cpu_graphik}

\begin{definition}[CPU-Graph]
\label{def:cpu_graph}
Die Architektur einer CPU wird modelliert als ein großer gewichteter gerichteter Multigraph $\Graph_{\mathrm{CPU}} = (V_{\mathrm{CPU}}, E_{\mathrm{CPU}}, w)$, wobei:
\begin{enumerate}[label=(\alph*)]
  \item Knoten $V_{\mathrm{CPU}}$ repräsentieren Register, ALU-Gatter, Speicherblatten.
  \item Kanten $E_{\mathrm{CPU}}$ repräsentieren Datenpfade und Signalleitungen.
  \item Kantengewichte $w(e)$ repräsentieren Signalwerte, Datenbreite, Verzögerungen.
\end{enumerate}
\end{definition}

\begin{remark}[Hierarchische Graphkomposition]
Eine praktische CPU ist nicht eine Zwei-Knoten-Graphen-Komposition, sondern eine hierarchische Komposition:
\[
  \Graph_{\mathrm{CPU}} = \bigcirc_{i=1}^{N} \Graph_i
\]
wobei jeder $\Graph_i$ ein Gatter oder ein Modul aus Gattern ist.

Diese hierarchische Struktur ist möglich, weil die Graphkomposition assoziativ ist (Theorem~\ref{thm:halbgruppe_halbleiter}).
\end{remark}

\subsection{Speicher und Zustände als Graphsche Schleifen}
\label{subsec:speicher_schleifen}

\begin{definition}[Speicher-Graph]
\label{def:speicher_graph}
Ein Speicherelement (Register, Flip-Flop) wird modelliert als ein Graph mit einer Schleife:
\[
  \Graph_{\mathrm{FF}} = (V_{\mathrm{FF}}, E_{\mathrm{FF}})
\]
mit:
\begin{enumerate}[label=(\alph*)]
  \item $V_{\mathrm{FF}} = \{v_{\mathrm{stored}}, v_{\mathrm{feedback}}\}$ (zwei Knoten).
  \item $E_{\mathrm{FF}} \ni (v_{\mathrm{stored}}, v_{\mathrm{feedback}})$ (Feedback-Kante).
  \item Kantengewicht $w(v_{\mathrm{stored}}, v_{\mathrm{feedback}}) = b$ (gespeichertes Bit).
\end{enumerate}
\end{definition}

\begin{theorem}[Stabilität von Speicher-Graphen]
\label{thm:speicher_stabilität}
Ein Speicher-Graph mit Feedback-Schleife ist stabil und speichert einen Zustand über mehrere Takte, solange die Feedback-Schleife erhalten bleibt.
\end{theorem}

\begin{proof}
Die Feedback-Schleife erzeugt eine positive Rückkopplung, die den Zustand verstärkt. Formal ist die Dynamik:
\[
  v_{\mathrm{stored}}^{(t+1)} = f(v_{\mathrm{stored}}^{(t)}, v_{\mathrm{input}}^{(t)})
\]

Wenn $v_{\mathrm{feedback}}$ direkt $v_{\mathrm{stored}}$ speist, wird der Zustand über Zeit erhalten.
\end{proof}

% ============================================================
\section{Zusammenfassung und Ausblick}
\label{sec:ausblick}
% ============================================================

\subsection{Zentrale Erkenntnisse}
\label{subsec:erkenntnisse}

Diese Arbeit hat eine fundamentale algebraische Grundlegung der Informatik und Elektronik durch Halbleiter-Graphen etabliert:

\begin{enumerate}[label=(\roman*)]
  \item \textbf{Binarität ist Minimum:} Der zweipolige Halbleiter mit zwei Zuständen (leitet/leitet nicht) ist das informationsminimalste elektronische Bauteil und kann genau ein Bit speichern.

  \item \textbf{Graphische Äquivalenz:} Jeder Halbleiter entspricht bijektiv einem Zwei-Knoten-Graphen. Diese Äquivalenz ist strukturerhaltend.

  \item \textbf{Komposition ist Algebra:} Die Komposition von Halbleiter-Graphen bildet eine Halbgruppe unter Matrizenmultiplikation (gemäß \cite{epp2026matrgraph}).

  \item \textbf{Gatter sind Kompositionen:} Alle logischen Gatter (AND, OR, NOT, etc.) werden durch Kompositionen von Halbleiter-Graphen realisiert.

  \item \textbf{Algorithmen sind Graphen:} Jeder Algorithmus und jede Berechnung ist äquivalent zu einer Graphkomposition von Halbleitern. Dies ist die fundamentale Natur der Informatik.

  \item \textbf{CPU ist Halbleiter-Netzwerk:} Eine CPU ist hierarchisch eine Komposition von Halbleiter-Graphen in vielen Schichten.
\end{enumerate}

\subsection{Philosophische Bedeutung}
\label{subsec:philosophie}

Der Halbleiter ist in der Tat der Schlüssel zum Leben. Nicht im metaphorischen Sinne, sondern im fundamentalsten mathematischen und physikalischen Sinne:

\begin{itemize}
  \item \textbf{Unentbehrlichkeit:} Ohne den Halbleiter gäbe es keine digitale Elektronik, keinen Computer, keine künstliche Intelligenz.

  \item \textbf{Eleganz der Binarität:} Die geniale Einfachheit von zwei Zuständen ermöglicht die ganze Komplexität der modernen Informatik.

  \item \textbf{Graphische Universalität:} Die graphische Darstellung durch Zwei-Knoten-Graphen zeigt, dass alle Informatik im Grunde eine Geometrie und Topologie ist.

  \item \textbf{Minimale Komplexität:} Der Halbleiter ist nicht nur minimal informativen Sinne, sondern auch das eleganteste Bauteil für Informationsspeicherung.
\end{itemize}

\subsection{Offene Forschungsfragen}
\label{subsec:offene_fragen}

Zukünftige Forschung könnte sich mit folgenden Fragen beschäftigen:

\begin{enumerate}[label=(\alph*)]
  \item \textbf{Quantenhalbleiter:} Wie lässt sich die graphische Algebra auf Quantenhalbleiter und Qubits erweitern?

  \item \textbf{Fehlerkorrektur:} Wie können Fehler in Halbleiter-Graphen graphisch modelliert und korrigiert werden?

  \item \textbf{Optische Halbleiter:} Können optische Transistoren (mit Photonen statt Elektronen) ebenfalls durch diese Graphenalgebra modelliert werden?

  \item \textbf{Neuromorphe Architekturen:} Wie können biologische Neuronen (die auch zwei Zustände: aktiv/inaktiv haben) durch diese Algebra modelliert werden?

  \item \textbf{Energieoptimierung:} Können graphische Methoden neue Wege zur Optimierung von Halbleiter-Schaltkreisen und Energieverbrauch öffnen?

  \item \textbf{Universelle Berechenbarkeit:} Wie könnte man umgekehrt zeigen, dass JEDE universell berechenbare Funktion notwendigerweise in dieser Halbleiter-Graphalgebra realisiert werden muss?

  \item \textbf{Biologische Information:} Sind biologische Informationsprozesse (DNA, Proteinsynthese) auch auf dieser fundamentalen Binarität aufgebaut?
\end{enumerate}

\subsection{Fazit}
\label{subsec:fazit}

Der Halbleiter ist nicht nur ein nützliches elektronisches Bauteil. Er ist die fundamentale Einheit der Informatik, der digitalen Information und letztlich des modernen Lebens. Durch die graphische Modellierung als Zwei-Knoten-Graphen haben wir eine saubere, algebraische Grundlegung geschaffen, die zeigt:

\begin{center}
\textit{Informatik ist die Komposition von Halbleiter-Graphen.}
\end{center}

Diese Sicht vereinheitlicht alle Bereiche: von der Elektronik über logische Gatter, Prozessorarchitektur bis hin zu Algorithmen und Turing-Berechenbarkeit. Sie zeigt auch, warum der Halbleiter erfolgreich geblieben ist: weil er die Minimale und eleganteste Lösung zum Problem der Informationsspeicherung darstellt.

In diesem Sinne ist der Halbleiter wirklich der Schlüssel zum Leben -- zum digitalen Leben, das unsere Welt heute prägt.

% ============================================================
\appendix
% ============================================================

\section{Ausführliche Beweise und technische Details}
\label{sec:anhang_beweise}

\subsection{Beweis: Äquivalenz von Halbleiter-Graphen und booleschen Funktionen}
\label{subsec:anhang_boolesche_aequivalenz}

\begin{theorem}[Bijektivität zwischen Zwei-Knoten-Graphen und 2×2-Matrizen]
\label{thm:anhang_2x2_bijektiv}
Es gibt eine Bijektion zwischen der Menge aller gewichteten gerichteten Zwei-Knoten-Graphen $\mathcal{G}_2$ und der Menge aller $2 \times 2$ reellen Matrizen $\mathcal{M}_{2 \times 2}(\R)$.
\end{theorem}

\begin{proof}
Die Bijektion wird gegeben durch die Adjazenzmatrix-Konstruktion:

Für einen Zwei-Knoten-Graph $\Graph = (\{v_1, v_2\}, E, w)$ ist die Adjazenzmatrix:
\[
  A_{\Graph} = \begin{pmatrix} w(v_1, v_1) & w(v_1, v_2) \\ w(v_2, v_1) & w(v_2, v_2) \end{pmatrix}
\]

Dies ist eine $2 \times 2$ Matrix. Die Umkehrabbildung ist trivial: Jede $2 \times 2$ Matrix definiert einen gewichteten Zwei-Knoten-Graphen durch ihre Einträge als Kantengewichte.
\end{proof}

\subsection{Vollständiger Beweis der funktionalen Vollständigkeit}
\label{subsec:anhang_vollstaendige_vollstaendigkeit}

\begin{lemma}[NAND ist funktional vollständig]
\label{lem:nand_vollstaendig}
Das NAND-Gatter $\uparrow: \Bool^2 \to \Bool$ mit $x \uparrow y = \neg(x \land y)$ ist funktional vollständig, d.h., jede boolesche Funktion kann aus NAND-Gattern konstruiert werden.
\end{lemma}

\begin{proof}
Es genügt zu zeigen, dass NOT und OR aus NAND konstruierbar sind:
\begin{enumerate}[label=(\roman*)]
  \item NOT: $\neg x = x \uparrow x$
  \item OR: $x \lor y = (x \uparrow x) \uparrow (y \uparrow y) = \neg x \uparrow \neg y = \neg(\neg x \land \neg y) = x \lor y$ (De Morgan)
\end{enumerate}

Da NAND funktional vollständig ist und NAND durch eine Halbleiter-Schaltung realisierbar ist (Negation eines AND-Gatters), ist die gesamte boolesche Logik durch Halbleiter realisierbar.
\end{proof}

\subsection{Detaillierte Analyse: CPU-Graphstruktur}
\label{subsec:anhang_cpu_detail}

\begin{definition}[Multi-Schicht-Graphkomposition]
\label{def:multi_schicht_komposition}
Eine hierarchische CPU wird modelliert als:
\[
  \Graph_{\mathrm{CPU}} = \bigcirc_{L=1}^{N} \bigcirc_{g=1}^{M_L} \Graph_{g}^{(L)}
\]
wobei:
\begin{enumerate}[label=(\alph*)]
  \item $L$ die Schicht-Tiefe ist (von Transistor-Ebene zu Gatter-Ebene bis zur Architekttur-Ebene).
  \item $M_L$ die Anzahl der Gatter/Module in Schicht $L$ ist.
  \item $\Graph_{g}^{(L)}$ der Graph des g-ten Elements in Schicht $L$ ist.
\end{enumerate}
\end{definition}

\begin{remark}[Hierarchische Abstraktion]
Diese hierarchische Struktur zeigt, dass:
\begin{itemize}
  \item Auf der untersten Ebene ($L=1$) sind alle Gates Zwei-Knoten-Graphen oder Kompositionen davon.
  \item Jede höhere Ebene ist eine Abstraktion (Komposition) der darunter liegenden Ebene.
  \item Am Ende ($L=N$) ist die ganze CPU selbst ein großer Zwei-Knoten-Graph (möglicherweise mit Millionen/Milliarden von Knoten).
\end{itemize}

Dies ist ein Beweis für die fundamentale Einfachheit und Einheit der Informatik.
\end{remark}

\subsection{Stabilitätsanalyse von Rückkopplungsschleifen}
\label{subsec:anhang_stabilitaet}

\begin{theorem}[Lyapunov-Stabilität für Speicher-Graphen]
\label{thm:anhang_lyapunov_speicher}
Ein Speicher-Graph mit positiver Feedback-Schleife ist Lyapunov-stabil für das Speichern eines binären Zustandes.
\end{theorem}

\begin{proof}
Sei $v(t)$ der Zustand des Speichers zum Zeitpunkt $t$. Die Rückkopplung ist:
\[
  v(t+1) = \sigma(w \cdot v(t) + u(t))
\]
wobei $\sigma$ eine Aktivierungsfunktion ist (z.B. Stufenfunktion), $w$ das Rückkopplungsgewicht und $u(t)$ die externe Eingabe.

Für $w > 1$ (positive Rückkopplung) und konstante Eingabe $u$ stabilisiert sich der Zustand auf einem festen Punkt $v^* = \sigma(w \cdot v^* + u)$.

Dies ist stabil im Sinne von Lyapunov, weil kleine Störungen $\delta v$ exponentiell abgebaut werden.
\end{proof}

% ============================================================
\section{Vergleich mit klassischen Informatik-Ansätzen}
\label{sec:klassische_ansaetze}

\subsection{Von Turing zur Graphkomposition}
\label{subsec:turing_zur_graph}

Die klassische Theorie der Berechenbarkeit basiert auf Turing-Maschinen \cite{turing1936}, die aus:
\begin{itemize}
  \item Einem unendlichen Bandspeicher,
  \item einem Lese-/Schreibkopf,
  \item einem Zustandsregister,
  \item einer Übergangsfunktion bestehen.
\end{itemize}

Unsere Graphkompositions-Sicht zeigt, dass diese abstrakte Struktur konkret durch Halbleiter realisierbar ist. Der Bandspeicher wird durch eine Folge von Speicher-Graphen (mit Feedback-Schleifen) realisiert. Der Zustand ist ein Knoten. Die Übergangsfunktion ist eine Graphkomposition von Halbleitern.

Dies ist nicht nur eine theoretische Äquivalenz, sondern eine praktische Realisierung.

\subsection{Von Boolean Logic zur Hardware}
\label{subsec:boolean_zu_hardware}

George Boole \cite{boole1854} entwickelte die Algebra der Logik rein mathematisch. Shannon \cite{shannon1938} zeigte, dass diese logik mit Relais realisiert werden konnte. Unsere Arbeit zeigt:

\begin{center}
\textit{Boolean Logic = Halbleiter-Graphkomposition = Berechenbarkeit}
\end{center}

Dies ist eine tiefe und überraschende Einheit.

\subsection{Von von Neumann zu Halbleiter-Architektur}
\label{subsec:von_neumann_halbleiter}

Von Neumanns Computerarchitektur \cite{neumann1945} besteht aus:
\begin{itemize}
  \item Control Unit (Steuerwerk),
  \item Arithmetic Logic Unit (Rechenwerk),
  \item Speicher,
  \item Ein- und Ausgabe.
\end{itemize}

Unsere Graphkompositions-Sicht zeigt, dass jede dieser Komponenten aus Halbleiter-Graphen aufgebaut ist und dass die Architektur selbst ein Halbleiter-Graphen-System ist.

% ============================================================
\bibliographystyle{plainnat}
\begin{thebibliography}{99}

\bibitem[Axler(2015)]{axler2015}
Axler, S. (2015).
\newblock \textit{Linear Algebra Done Right}, 3rd ed.
\newblock Springer.

\bibitem[Boole(1854)]{boole1854}
Boole, G. (1854).
\newblock \textit{An Investigation of the Laws of Thought on Which Are Founded the Mathematical Theories of Logic and Probabilities}.
\newblock Macmillan.

\bibitem[Diestel(2017)]{diestel2017}
Diestel, R. (2017).
\newblock \textit{Graph Theory}, 5th ed.
\newblock Springer.

\bibitem[Epp(2026)]{epp2026matrgraph}
Epp, S. (2026).
\newblock \textit{Äquivalenz durch Sicht: Graphische Deduktion und die Überflüssigkeit der Matrix}.
\newblock September 2026.

\bibitem[Godsil and Royle(2001)]{godsil2001}
Godsil, C. and Royle, G. (2001).
\newblock \textit{Algebraic Graph Theory}.
\newblock Springer.

\bibitem[Horn and Johnson(2013)]{horn2013}
Horn, R.~A. and Johnson, C.~R. (2013).
\newblock \textit{Matrix Analysis}, 2nd ed.
\newblock Cambridge University Press.

\bibitem[Khalil(2002)]{khalil2002}
Khalil, H.~K. (2002).
\newblock \textit{Nonlinear Systems}, 3rd ed.
\newblock Prentice Hall.

\bibitem[Neumann(1945)]{neumann1945}
von Neumann, J. (1945).
\newblock First Draft of a Report on the EDVAC.
\newblock \textit{IEEE Annals of the History of Computing}, 15(4):27--75.

\bibitem[Newman(2010)]{newman2010}
Newman, M. (2010).
\newblock \textit{Networks: An Introduction}.
\newblock Oxford University Press.

\bibitem[Shannon(1938)]{shannon1938}
Shannon, C.~E. (1938).
\newblock A Symbolic Analysis of Relay and Switching Circuits.
\newblock \textit{Transactions of the American Institute of Electrical Engineers}, 57(12):713--723.

\bibitem[Spielman(2019)]{spielman2019}
Spielman, D.~A. (2019).
\newblock \textit{Spectral and Algebraic Graph Theory}.
\newblock Yale University.

\bibitem[Turing(1936)]{turing1936}
Turing, A.~M. (1936).
\newblock On Computable Numbers, with an Application to the Entscheidungsproblem.
\newblock \textit{Proceedings of the London Mathematical Society}, 42(1):230--265.

\bibitem[West(2001)]{west2001}
West, D.~B. (2001).
\newblock \textit{Introduction to Graph Theory}, 2nd ed.
\newblock Prentice Hall.

\end{thebibliography}

\end{document}